\scanpage{345}
% top centre
\mspage{261}
% lines 1–7
Il est temps enfin de faire appel à l'hypothèse c) de 4.10.3. Soit donc $\alpha$ un ordinal limite, comme stipulé dans cette condition. Comme $\mathrm{card}\,\alpha\leq\pi$, \unsure{à|a} fortiori $\mathrm{card}\,\beta\leq\pi$ si $\beta<\alpha$, on a donc un système inductif
% line 8
\[
 \bigl(X_\beta(\cdot)\bigr)_{\beta\in\unsure{I|J}_\alpha}
\]
% lines 9–11
dans $\prod_{j\in J}\unsure{\mathcal E|\mathscr E|\mathcal C}_j$. Comme les $\varinjlim$ de type $I_\alpha$ \unsure{existent|existant} dans $\mathcal E_j$, on \unsure{peut} \deleted{\unsure{co}} donc considérer
% line 12
\[
 G(j)=\varinjlim_{\beta\in I_\alpha}X_\beta(j)\in\operatorname{Ob}\mathcal E_j
\]

% lines 13–14
\unsure{Si} $\lambda:j'\unsure{\mathrel{-}|\longrightarrow}j$ est une flèche de $J$, on en déduit une $\lambda$-flèche\unsure{s|}
% line 15, above display
\deleted{$G(\lambda):G(j)=\unsure{l}$}
% line 15
\[
 G(\lambda):G(j)=\varinjlim X_\beta(j)\longrightarrow
 G(j')\unsure{\simeq|=}\varinjlim X_{\unsure{\beta+1|\beta}}(j')
\]
% line 16
par $G(\lambda)=\varinjlim t_\beta(\lambda)$. Il faut
% archive sheet number: 344
